% tenkz-core — L0: semantic hues, theme slots, styles, the metric system,
% the label-placement constants, and the event stream.
%
% Metric doctrine: ONE base pitch; every repeated distance is a named ratio
% of it.  Profiles (compact, inline) are a scale factor applied at the base,
% so no literal can bypass them.  Each ratio's motivation is printed in the
% manual's geometry appendix; the source states it once here.
%
% Label doctrine: labels never overlap ink by construction.  Every label
% site (in-glyph, leg tip, bond midpoint, free quadrant of a bead) is a
% reserved band derived from the metric system; automatic placement picks a
% free quadrant, and every label accepts `label pos` / `label shift` for
% manual override.
%
% Tensor-style doctrine: ONE glyph policy per scope decides how an
% unadorned tensor (\tn{A}, \tnput{...}) renders.
%   tensor style=dot   (default) filled bead, label placed OUTSIDE in a
%                      free quadrant by the auto-quadrant algorithm
%                      (the CPSV school);
%   tensor style=box   square/rectangular box, label typeset INSIDE
%                      (the Bultinck/Ghent school); \tn*{A} inscribes
%                      \overline{A}, and no external auto label exists,
%                      so brace clearances tighten accordingly.
% The policy is a document-wide \tnset key and a per-picture key of the
% tenkz environment, \tnpic, and tenkzfree (the environments group, so a
% picture-local assignment scopes itself; tenkzcd cells inherit it through
% the \tnpic objects they contain).  Explicit per-cell skins (dot, box, pill,
% mpo) always override the policy; the event stream always records the
% RESOLVED skin.

% Triangle skins (tri=l / tri=r) borrow the isosceles-triangle node shape.
\usetikzlibrary{shapes.geometric}

% ---------- semantic hues (a theme may rebind colours, nothing else) ----------
\colorlet{tenkzInk}{black!92}
\colorlet{tenkzPaper}{white}
\colorlet{tenkzPassive}{black!55}
\colorlet{tenkzAction}{red!65!black}
\colorlet{tenkzMarked}{blue!65!black}
\colorlet{tenkzExtra}{violet!65!black}
\colorlet{tenkzOperator}{red!80!black}% wire hue of an operator (MPO) layer

% ---------- metric system ----------
\newdimen\tenkz@pitch          \tenkz@pitch=11mm
\newdimen\tenkz@basepitch      \tenkz@basepitch=11mm
% ratios (of pitch unless stated)
\def\tenkz@r@layersep{0.80}    % layers read as layers, not as two chains
\def\tenkz@r@physleg{0.38}     % a labelled leg clears the bond-label band
\def\tenkz@r@stub{0.45}        % open index: longer than a leg, shorter than a bond;
                               % also the reach of an inter-layer cup: the bend
                               % protrudes exactly as far as the stubs it replaces,
                               % so open and cup-closed ends share one silhouette
\def\tenkz@r@traceclear{0.62}  % trace racetrack clears leg labels on the outer row
\def\tenkz@r@labelclear{0.12}  % ONE clearance for every label band
\def\tenkz@r@dotlabelband{0.30}% depth of a script-size dot label: clearance + text
\def\tenkz@r@inlinephysleg{0.26}   % inline: legs shrink faster than the pitch
\def\tenkz@r@inlinelabelclear{0.05}% inline: labels hug the ink to spare line height
\def\tenkz@r@dotdia{0.16}      % bead diameter
\def\tenkz@r@tallover{0.34}    % vertical overshoot of a wires=k glyph beyond its
                               % outer wire rows: the gate box must visibly COVER
                               % the wires it terminates (quantikz \gate[k])
\def\tenkz@r@phtracereach{0.24}% physical-trace clearance: how far east of the
                               % glyph edge the per-site Tr_phys loop swings; it
                               % must clear the glyph yet stay inside the bond
                               % span to the next site
\def\tenkz@r@pairtracereach{0.40}% pair-trace reach (traced columns): how far east
                               % of the leg axis the per-column Tr_s racetrack
                               % swings; < 1/2 so the east straight clears the
                               % neighbouring column's glyph and label bands at
                               % the default pitch, and wide enough that the
                               % stadium caps (radius = reach/2) read as smooth
                               % turnarounds rather than kinks
\def\tenkz@r@regionmargin{0.36}% < pitch/2 so notches read; > glyph radius
\def\tenkz@r@corner{0.14}      % rounded-corner radius of traces and hulls
% absolute floors and ink hierarchy
\def\tenkz@wirewidth{0.55pt}
\def\tenkz@annwidth{0.40pt}
\def\tenkz@emphwidth{0.80pt}
\def\tenkz@junctionfloor{0.9mm}% below ~3x wire width a junction reads as dirt

\def\tenkz@dim#1{\dimexpr#1\tenkz@pitch\relax}

% ---------- glyph policy (see the tensor-style doctrine above) ----------
% \tenkz@tensorstyle holds the resolved mode, `dot` or `box`; grid cells
% and \tnput read it as their default skin at option-parse time.
\def\tenkz@tensorstyle{dot}

% ---------- pgfkeys front door ----------
\def\tnset#1{\pgfqkeys{/tenkz}{#1}}
\pgfkeys{
  /tenkz/.is family,
  /tenkz/pitch/.code={\tenkz@pitch=#1\relax},
  /tenkz/compact/.code={\tenkz@pitch=0.8\tenkz@basepitch},
  /tenkz/inline/.code={\tenkz@pitch=0.62\tenkz@basepitch
                       \let\tenkz@labelsize\scriptscriptstyle
                       \let\tenkz@r@physleg\tenkz@r@inlinephysleg
                       \let\tenkz@r@labelclear\tenkz@r@inlinelabelclear},
  /tenkz/tensor style/.is choice,
  /tenkz/tensor style/dot/.code={\def\tenkz@tensorstyle{dot}},
  /tenkz/tensor style/box/.code={\def\tenkz@tensorstyle{box}},
}

% Label mathematics is script-size by default: glyphs are compact marks, not
% text boxes, and the surrounding equation stays visually dominant.
\let\tenkz@labelsize\scriptstyle

% Every inscribed glyph label carries this strut, so boxes in one row have
% one height whether their labels are A, T_g, or A^{(2)} (Tufte: small
% multiples must be uniform; a descender is not information).
\def\tenkz@glyphstrut{\vphantom{A^{X}_{g}}}

% ---------- style table (complete Phase-0 subset of the spec table) ----------
\tikzset{
  tenkz every picture/.style={line cap=round, line join=round},
  % glyph skins
  tensor/.style={draw=tenkzInk, fill=tenkzInk, circle, inner sep=0pt,
      minimum size=\tenkz@dim{\tenkz@r@dotdia}, line width=\tenkz@wirewidth},
  conjugate tensor/.style={tensor},
  box tensor/.style={draw=tenkzInk, fill=tenkzPaper, rectangle,
      inner sep=1.6pt, minimum size=\tenkz@dim{0.22},
      line width=\tenkz@wirewidth, text=tenkzInk},
  mpo tensor/.style={box tensor},
  pill tensor/.style={draw=tenkzInk, fill=tenkzPaper, rectangle,
      rounded corners=2.2pt, inner xsep=3.2pt, inner ysep=2.2pt,
      minimum height=\tenkz@dim{0.22}, line width=\tenkz@wirewidth,
      text=tenkzInk},
  % capsule, not circle: the height is pinned by `minimum size`, and a wide
  % label (an inverse, a product) grows the glyph only sideways, so every
  % on-wire matrix on a chain reads at the same height
  on-wire matrix/.style={draw=tenkzAction, fill=tenkzPaper, rectangle,
      rounded corners=\tenkz@dim{0.15}, inner xsep=1.8pt, inner ysep=1.4pt,
      minimum size=\tenkz@dim{0.30},
      line width=\tenkz@wirewidth, text=tenkzInk},
  % the style's public 0.6 name; \tnX is the documented alias
  ring tensor/.style={on-wire matrix},
  % canonical-form isometries (tri=l / tri=r): an isoceles triangle whose FLAT
  % side receives the isometry's domain and whose APEX emits its codomain --
  % the A_L / A_R glyphs of canonical-form MPS (sum_s A_L^s{}^dagger A_L^s = 1).
  % Wires meet the flat side's midpoint and the apex, both on the wire axis.
  % inner sep is tight (0.6pt): a triangle must CIRCUMSCRIBE its text box,
  % so every point of padding costs ~triple in glyph height; the inscribed
  % label also drops the glyph strut (see \tenkz_tri_node:nnn) or two
  % triangle rows would collide across the default layer gap
  % apex angle 55: the flat-side height grows like 2 tan(apex/2) * text
  % width, so a narrower apex trades a slightly longer glyph for the
  % vertical clearance the sandwich genre needs across the layer gap
  canonical tensor/.style={draw=tenkzInk, fill=tenkzPaper,
      isosceles triangle, isosceles triangle apex angle=55,
      isosceles triangle stretches, inner sep=0.6pt,
      minimum width=\tenkz@dim{0.34}, minimum height=\tenkz@dim{0.30},
      line width=\tenkz@wirewidth, text=tenkzInk},
  left canonical tensor/.style={canonical tensor, shape border rotate=0},
  right canonical tensor/.style={canonical tensor, shape border rotate=180},
  fusion map/.style={draw=tenkzInk, fill=tenkzInk, line cap=round,
      line width=1.4pt},
  tree junction/.style={draw=tenkzInk, fill=tenkzInk, circle, inner sep=0pt,
      minimum size=\tenkz@junctionfloor},
  % wires
  bond/.style={draw=tenkzInk, line width=\tenkz@wirewidth},
  fused bond/.style={bond, double=tenkzPaper, double distance=1.1pt},
  physical leg/.style={bond},
  operator bond/.style={bond, draw=tenkzOperator},
  trace/.style={bond, rounded corners=\tenkz@dim{\tenkz@r@corner}},
  % ---------- semantic roles (the 0.6 hue contract) ----------
  % Resolution is pure style indirection: role words map to ONE derived
  % style per role x glyph family — bond (wires, joins), tensor (filled
  % beads), outline (inscribed-label glyphs: box, pill, mpo, tri, ring),
  % edge and site (the lattice families) — all bound to the core hue
  % slots, so a theme that rebinds tenkzOperator etc. recolours every
  % role-hued atom, wire and join at one point.  No raw colour is ever
  % accepted at a call site, in any tier.
  marked bond/.style={bond, draw=tenkzMarked},
  extra bond/.style={bond, draw=tenkzExtra},
  passive bond/.style={bond, draw=tenkzPassive},
  operator tensor/.style={draw=tenkzOperator, fill=tenkzOperator},
  marked tensor/.style={draw=tenkzMarked, fill=tenkzMarked},
  extra tensor/.style={draw=tenkzExtra, fill=tenkzExtra},
  passive tensor/.style={draw=tenkzPassive, fill=tenkzPassive},
  operator outline/.style={draw=tenkzOperator, text=tenkzOperator},
  marked outline/.style={draw=tenkzMarked, text=tenkzMarked},
  extra outline/.style={draw=tenkzExtra, text=tenkzExtra},
  passive outline/.style={draw=tenkzPassive, text=tenkzPassive},
  operator edge/.style={distinguished edge, draw=tenkzOperator},
  marked edge/.style={distinguished edge, draw=tenkzMarked},
  extra edge/.style={distinguished edge, draw=tenkzExtra},
  passive edge/.style={distinguished edge, draw=tenkzPassive},
  operator site/.style={draw=tenkzOperator, fill=tenkzOperator},
  marked site/.style={draw=tenkzMarked, fill=tenkzMarked},
  extra site/.style={draw=tenkzExtra, fill=tenkzExtra},
  passive site/.style={draw=tenkzPassive, fill=tenkzPassive},
  % directionality marks (design brief): an OUTGOING arrow is the vector space
  % V, an INCOMING one its dual V*; contraction joins out-to-in, and REVERSING
  % a barb is the dual/inverse representation -- so `reversed` variants flip
  % the barb, never the wire.  The barb rides mid-wire and inherits the wire's
  % colour; the `fused` variants size the barb to span the doubled wire.
  dir marks/.style={decoration={markings, mark=at position 0.55 with
      {\arrow{Straight Barb[length=2.6pt]}}}, postaction=decorate},
  dir marks reversed/.style={decoration={markings, mark=at position 0.55 with
      {\arrow{Straight Barb[length=2.6pt, reversed]}}}, postaction=decorate},
  fused dir marks/.style={decoration={markings, mark=at position 0.55 with
      {\arrow{Straight Barb[length=3.8pt]}}}, postaction=decorate},
  fused dir marks reversed/.style={decoration={markings,
      mark=at position 0.55 with
      {\arrow{Straight Barb[length=3.8pt, reversed]}}}, postaction=decorate},
  bond arrow/.style={bond, dir marks},
  fused bond arrow/.style={fused bond, fused dir marks},
  % annotation ink (never confusable with 0.55pt wires)
  cut/.style={draw=tenkzPassive, dashed, line width=\tenkz@annwidth},
  brace/.style={draw=tenkzInk, line width=\tenkz@annwidth,
      decoration={brace, amplitude=2.6pt}, decorate},
  label/.style={text=tenkzInk, inner sep=1pt, font=\scriptsize},
  cd arrow/.style={draw=tenkzInk, line width=\tenkz@annwidth,
      -{Straight Barb[length=3.4pt]},
      shorten <=3.5pt, shorten >=3.5pt},
  % regions
  region selected/.style={draw=tenkzAction, fill=tenkzAction!12!tenkzPaper,
      line width=\tenkz@emphwidth},
  region secondary/.style={draw=tenkzMarked, fill=tenkzMarked!10!tenkzPaper,
      line width=\tenkz@emphwidth},
  region complement/.style={draw=tenkzPassive,
      fill=tenkzPassive!8!tenkzPaper, densely dashed,
      line width=\tenkz@emphwidth},
  region collar/.style={draw=tenkzExtra, fill=tenkzExtra!10!tenkzPaper,
      line width=\tenkz@emphwidth},
  region outline/.style={fill=none},
  distinguished edge/.style={draw=tenkzAction, line width=\tenkz@emphwidth},
  removed site/.style={draw=tenkzPassive, fill=tenkzPaper,
      line width=\tenkz@wirewidth},
}

% ---------- event stream (v2, minimal Phase-0 writer) ----------
\newwrite\tenkz@log
\newif\iftenkz@logopen
% Two registers, two jobs.  \tenkz@pictureuid is the GLOBAL uniqueness
% counter: it only ever grows, so every picture's id is fresh.
% \tenkz@pictureid is the EMITTING context: it is assigned LOCALLY at each
% picture's begin (every begin site sits inside its picture's group), so
% closing the group RESTORES the enclosing picture's id -- a nested \tnpic
% inside a tenkzcd cell must not steal the parent's later events (cdcell
% etc.), and back in running prose the register reverts to 0, which is how
% a bare \tntree logs picture=0.
\newcount\tenkz@pictureid
\newcount\tenkz@pictureuid
\AtBeginDocument{\immediate\openout\tenkz@log=\jobname.tnlog\tenkz@logopentrue}
\AtEndDocument{\iftenkz@logopen\immediate\closeout\tenkz@log\fi}
\def\tenkz@event#1{\iftenkz@logopen
  \immediate\write\tenkz@log{#1}\fi}
\def\tenkz@beginpicture#1{\global\advance\tenkz@pictureuid by 1\relax
  \tenkz@pictureid=\tenkz@pictureuid
  \tenkz@event{picture|id=\the\tenkz@pictureid|lang=#1}}

% ---------- species (the 0.6 hue contract, clause 2) ----------
% Species declare NAMES, never colors: `species={right,left}' (document
% preamble via \tnset, or picture scope) assigns each new name the next
% hue of the theme's semantic cycle and derives one style per glyph
% family — `species <name> bond', `species <name> tensor', `species
% <name> outline' — through the same indirection as the role styles, so
% a theme rebinding the core hue slots recolours every species at one
% point.  A name keeps its index (hence its hue) across redeclarations,
% so a document-global declaration holds one hue per species over every
% figure of a paper.  An undeclared species at a call site is an error.
\ExplSyntaxOn
\prop_new:N \g__tenkz_species_prop     % name -> 1-based declaration index
\int_new:N  \g__tenkz_species_int
% the semantic cycle: hue slots, in assignment order
\clist_const:Nn \c__tenkz_speccycle_clist
  { tenkzMarked , tenkzAction , tenkzExtra , tenkzOperator }
\msg_new:nnn { tenkz } { unknown-species }
  {
    Undeclared~species~'#1'.~Declare~names~first~--~
    \token_to_str:N \tnset{species={...}}~in~the~preamble~(recommended:~
    a~species~keeps~one~hue~across~every~figure)~or~in~the~picture's~
    option~list.
  }
\cs_new_protected:Npn \tenkz_species_check:n #1
  {
    \prop_if_in:NnF \g__tenkz_species_prop {#1}
      { \msg_error:nnn { tenkz } { unknown-species } {#1} }
  }
\cs_generate_variant:Nn \tenkz_species_check:n { e, V }
\cs_new_protected:Npn \tenkz_species_declare:n #1
  { \clist_map_inline:nn {#1} { \tenkz_species_one:n {##1} } }
\cs_new_protected:Npn \tenkz_species_one:n #1
  {
    % first declaration fixes the index; a redeclaration re-derives the
    % styles (idempotent) without moving any later species' hue
    \prop_if_in:NnF \g__tenkz_species_prop {#1}
      {
        \int_gincr:N \g__tenkz_species_int
        \prop_gput:Nne \g__tenkz_species_prop {#1}
          { \int_use:N \g__tenkz_species_int }
      }
    \prop_get:NnN \g__tenkz_species_prop {#1} \l_tmpa_tl
    \tl_set:Ne \l_tmpb_tl
      { \clist_item:Nn \c__tenkz_speccycle_clist
          { 1 + \int_mod:nn { \l_tmpa_tl - 1 }
              { \clist_count:N \c__tenkz_speccycle_clist } } }
    \tl_set:Ne \l_tmpa_tl
      {
        species~#1~bond/.style={bond, draw=\l_tmpb_tl},
        species~#1~tensor/.style={draw=\l_tmpb_tl, fill=\l_tmpb_tl},
        species~#1~outline/.style={draw=\l_tmpb_tl, text=\l_tmpb_tl}
      }
    \exp_args:NV \tikzset \l_tmpa_tl
  }
\pgfkeys{
  /tenkz/species/.code={\tenkz_species_declare:n{#1}},
}
\ExplSyntaxOff

\endinput
